\begin{pastexam}[label={px:ncku-114-cce-ds-03}]{成大電通所 114 學年度 單選題第 3 題}
（單選；每題 5 分，答錯倒扣 2 分）

Given a directed graph in the following, please identify which edge would become a cross edge if we traverse
the graph by the breadth first search algorithm from node C. Assume that the nodes are stored in the adjacent
list according to alphabetic order. Please specify which edge forms a cross edge in the resulting breadth-first
spanning tree.

\begin{center}
\begin{tikzpicture}[v/.style={circle,draw,inner sep=1pt,minimum size=7.5mm},>={Stealth[length=2.4mm]}]
\node[v] (A) at (0,0) {A};
\node[v] (B) at (1.9,1.15) {B};
\node[v] (C) at (2.4,-0.8) {C};
\node[v] (D) at (4.9,0) {D};
\node[v] (F) at (7.0,1.6) {F};
\node[v] (E) at (7.0,0) {E};
\draw[->] (A) -- (B);
\draw[->] (A.-40) -- (C.195);
\draw[->] (C.160) -- (A.-5);
\draw[->] (B.-95) -- (C.130);
\draw[->] (C.95) -- (B.-60);
\draw[->] (B.-30) -- (D.170);
\draw[->] (D.135) -- (B.5);
\draw[->] (C) -- (D);
\draw[->] (D) -- (F);
\draw[->] (D) -- (E);
\draw[->] (F) -- (E);
\end{tikzpicture}
\end{center}
\begin{enumerate}[label=\alph*),nosep]
\item (F, E)
\item (B, D)
\item (D, B)
\item (A, B)
\item All of the above.
\end{enumerate}
\end{pastexam}

\begin{pastsoln}
答案：e) All of the above。

圖的邊：A$\to$B、A$\to$C、B$\to$C、B$\to$D、C$\to$A、C$\to$B、C$\to$D、D$\to$B、D$\to$E、D$\to$F、F$\to$E。
鄰接串列依字母順序，從 C 做 BFS（資結書第~21.3~節）：
\begin{center}
\begin{tabular}{lll}
\toprule
取出 & 依序看的邊 & 佇列（取出後） \\
\midrule
C（層 0） & C$\to$A、C$\to$B、C$\to$D 都是新點，樹邊 & A, B, D \\
A（層 1） & A$\to$B（B 已發現）、A$\to$C（C 已發現） & B, D \\
B（層 1） & B$\to$C（已發現）、B$\to$D（已發現） & D \\
D（層 1） & D$\to$B（已發現）、D$\to$E、D$\to$F 新點，樹邊 & E, F \\
E（層 2） & 沒有出邊 & F \\
F（層 2） & F$\to$E（已發現） & 空 \\
\bottomrule
\end{tabular}
\end{center}
BFS 樹：C 是根，子節點 A、B、D；E、F 是 D 的子節點。

非樹邊依兩端在樹上的關係分類（資結書第~21.5~節）：終點是起點的祖先叫後向邊（back edge），兩端沒有祖孫關係叫交叉邊
（cross edge）。 有向圖的 BFS 不會有前向邊。
\begin{itemize}
\item (F, E)：F、E 都在層 2，是兄弟（都是 D 的子節點），沒有祖孫關係，交叉邊。
\item (B, D)：B、D 都在層 1，都是 C 的子節點，交叉邊。
\item (D, B)：同上，反方向，也是交叉邊。
\item (A, B)：A、B 都在層 1，交叉邊。
\end{itemize}
四條都是交叉邊，選 e)。 其餘的非樹邊 A$\to$C、B$\to$C 指向根 C，是後向邊。
\end{pastsoln}
